\documentclass[11pt,a4paper]{article}

\usepackage{../shared/cathedral-preamble}

\title{%
  The Riemann Hypothesis and Cryptographic Security:\\
  What the Cathedral Proves, What It Doesn't,\\
  and What the Headlines Will Get Wrong%
}
\author{
  Jason Robert Gochanour
}
\date{April 20, 2026}

\begin{document}

\maketitle

\begin{abstract}
We analyze the cryptographic security implications of the Cathedral, a
machine-verified Lean~4 formalization establishing that the Riemann
Hypothesis is equivalent to the decay of the Nyman--Beurling distance.
We rigorously demonstrate that this result---and indeed any resolution
of the Riemann Hypothesis, whether proved or disproved---has
\textbf{no direct impact} on the security of RSA, elliptic curve
cryptography, AES, SHA-256, or any deployed cryptographic system.
We identify and debunk seven common misconceptions that security
researchers and journalists may propagate, provide precise
separations between distributional knowledge (which RH provides) and
computational hardness (which cryptography requires), and analyze
the one scenario in which RH-adjacent mathematics \emph{could}
theoretically affect cryptography: the discovery of a new structural
property of primes during the proof process. We conclude that the
Cathedral is cryptographically neutral, and that proving the Riemann
Hypothesis would, if anything, \emph{strengthen} the theoretical
foundations of existing cryptographic systems by converting
conditional results to unconditional ones.
\end{abstract}

\tableofcontents

% ================================================================
\section{Introduction: Why This Paper Is Necessary}
% ================================================================

When media outlets report on progress toward the Riemann Hypothesis,
the following headline pattern is nearly inevitable:

\begin{quote}
\emph{``Scientists Crack 167-Year-Old Math Problem --- Is Your
Encryption at Risk?''}
\end{quote}

The answer is no. But the reasons are subtle enough that they deserve
a rigorous treatment, particularly because the Cathedral involves
primes, zeta functions, and $L^2$ approximation theory --- all of
which sound superficially connected to cryptography.

This paper serves three audiences:

\begin{enumerate}
  \item \textbf{Security researchers} who need to assess whether the
    Cathedral warrants a threat advisory.
  \item \textbf{Journalists} who need accurate language for public
    reporting.
  \item \textbf{Cryptographers} who want precise statements about
    what distributional results (like RH) do and do not imply about
    computational hardness.
\end{enumerate}

\subsection{Summary of Findings}

\begin{center}
\fbox{\parbox{0.85\textwidth}{
\textbf{Bottom line:} The Cathedral (and any resolution of RH)
has zero impact on the security of:
\begin{itemize}[noitemsep]
  \item RSA (factoring-based)
  \item ECDSA / EdDSA (elliptic curve-based)
  \item AES-128/256 (symmetric encryption)
  \item SHA-2 / SHA-3 (hash functions)
  \item CRYSTALS-Kyber / Dilithium (lattice-based, post-quantum)
  \item Signal Protocol, TLS 1.3, or any deployed system
\end{itemize}
}}
\end{center}

% ================================================================
\section{What the Cathedral Actually Proves}
\label{sec:what}
% ================================================================

The crown theorem of the Cathedral states:
\begin{theorem}[Nyman--Beurling Equivalence, Machine-Verified]
\[
  \RH \;\iff\;
  \forall\,\varepsilon > 0,\;\exists\,N_0,\;\forall\,N \geq N_0,\;
  \exists\,v \in \RR^{N-1} :\quad
  \int_0^1 \biggl(1 - \sum_{k=2}^{N} v_k\,\fract{1/(kx)}\biggr)^2\,dx < \varepsilon.
\]
\end{theorem}

This is a statement about the density of a function space. It says
nothing about:
\begin{itemize}
  \item Factoring integers.
  \item Computing discrete logarithms.
  \item Finding collisions in hash functions.
  \item Solving lattice problems.
  \item Any algorithmic task used in cryptography.
\end{itemize}

The Cathedral is a \emph{structural equivalence}, not an algorithm.
It reduces one mathematical statement (RH) to another equivalent
mathematical statement ($d_N^2 \to 0$). Neither statement provides
a computational shortcut for any known hard problem.

% ================================================================
\section{Seven Myths Debunked}
\label{sec:myths}
% ================================================================

\begin{myth}[``Proving RH breaks RSA.'']
\label{myth:rsa}
\textbf{False.}
\end{myth}

RSA security relies on the computational hardness of factoring
$n = pq$ where $p$ and $q$ are large primes. Specifically:

\begin{definition}[Factoring Assumption]
There is no classical algorithm that, given a uniformly random
$n$-bit semiprime $N = PQ$, outputs $P$ (or $Q$) in time
polynomial in $n$, with non-negligible probability.
\end{definition}

The Riemann Hypothesis is a statement about the \emph{distribution}
of primes: it constrains the error term in the Prime Number Theorem
to $O(\sqrt{x}\log^2 x)$ rather than the unconditional
$O(x\exp(-c\sqrt{\log x}))$.

\begin{claim}[Distribution $\neq$ Factoring]
\label{claim:separation}
Knowing the exact distribution of primes does not help factor a
specific semiprime. The analogy: knowing that 50\% of the population
is female does not help you determine the gender of a specific
individual.
\end{claim}

More precisely: the prime-counting function $\pi(x)$ tells you
\emph{how many} primes exist below $x$. Factoring requires knowing
\emph{which} integers near $\sqrt{N}$ are prime and divide $N$.
These are fundamentally different computational problems. The first
is solved by RH (with optimal error bounds); the second is
conjectured to be hard regardless of RH.

\begin{remark}
Cryptographers already routinely \emph{assume} RH is true. Many
theoretical results in cryptography are proved ``conditional on GRH''
(the Generalized Riemann Hypothesis). If RH were secretly useful for
breaking RSA, someone would have already exploited the assumption.
\end{remark}

\begin{myth}[``Proving RH reveals a pattern in the primes that
could be exploited.'']
\label{myth:pattern}
\textbf{False.}
\end{myth}

RH constrains the \emph{global statistical distribution} of primes.
It tells you that primes of size $\sim x$ occur with density
$\sim 1/\ln x$, with fluctuations of size $\sim \sqrt{x}$. This is
already assumed in every cryptographic security analysis.

The ``pattern'' that RH reveals is:
\begin{quote}
\emph{Primes are distributed as regularly as possible, subject to
their arithmetic constraints.}
\end{quote}

This is the \emph{opposite} of a vulnerability. Irregularity in
prime distribution would be more concerning for cryptography than
regularity. If primes clustered unexpectedly, it might create
exploitable structure in cryptographic keys. RH asserts they
\emph{don't} cluster---they're as well-distributed as possible.

\begin{myth}[``The Möbius function in the proof could be used
to factor numbers.'']
\label{myth:mobius}
\textbf{False.}
\end{myth}

The Cathedral uses the Mertens function $M(x) = \sum_{n \leq x} \mu(n)$,
where $\mu(n)$ is the M\"obius function. Security researchers may note
that $\mu(n)$ encodes the factorization of $n$:
\[
  \mu(n) = \begin{cases}
    (-1)^k & \text{if } n \text{ is a product of } k \text{ distinct primes} \\
    0 & \text{if } n \text{ has a squared prime factor}
  \end{cases}
\]

The circularity is immediate: \textbf{computing $\mu(n)$ requires
knowing the factorization of $n$.} You cannot use $\mu(n)$ to factor
$n$---you need the factorization to compute $\mu(n)$ in the first
place.

The Mertens function $M(x)$ is a \emph{sum} over all $n \leq x$,
which encodes statistical information about the distribution of
prime factors across all integers, not about any specific integer.
It is computable in $O(x^{2/3})$ time without factoring anything
(via the Meissel--Mertens method), but this computation provides no
information useful for factoring a specific $N$.

\begin{myth}[``The Gram matrix eigenvalues could reveal prime
factors.'']
\label{myth:gram}
\textbf{False.}
\end{myth}

The Gram matrix $G_N$ has entries
$G(j,k) = \int_0^1 \fract{1/(jx)} \fract{1/(kx)}\,dx$, and its
eigenvalues depend on the arithmetic interaction between integers
$2 \leq j, k \leq N$. Security researchers might wonder: could the
eigenvalue structure of $G_N$ (for $N$ near a target semiprime)
reveal factoring information?

No. The Gram matrix depends on $\gcd(j,k)$ through the Vasyunin
formula, but this is the GCD of the \emph{row/column indices}, not
of the target number. The matrix $G_N$ is a fixed mathematical
object that depends only on $N$, not on any secret key or private
value. It is publicly computable. There is nothing to ``reveal.''

\begin{myth}[``The zeta zeros are like private keys---knowing them
breaks the system.'']
\label{myth:zeros}
\textbf{False.}
\end{myth}

The nontrivial zeros of $\zeta(s)$ are public mathematical constants.
The first $10^{13}$ zeros have been computed and published. Databases
of zeta zeros are freely available online (LMFDB, Odlyzko's tables).
Knowing billions of zeta zeros has not helped anyone factor any
specific integer.

The zeros encode the \emph{Fourier spectrum} of the prime counts,
not the location of individual primes. They are the frequency
content of the signal $\pi(x) - \text{Li}(x)$, which tells you
about the global oscillation of the prime-counting function, not
about specific values.

\begin{myth}[``Disproving RH would be worse for cryptography.'']
\label{myth:disproof}
\textbf{Also false} (with a nuance).
\end{myth}

If RH were \emph{disproved}---if a zero were found off the critical
line at $\rho = \sigma + it$ with $\sigma \neq 1/2$---the
consequences would be:

\begin{enumerate}
  \item The prime-counting function $\pi(x)$ would have larger
    fluctuations than currently assumed: $O(x^\sigma)$ instead
    of $O(x^{1/2})$.
  \item Certain conditional results in complexity theory (proved
    assuming GRH) would lose their theoretical backing.
  \item The Mertens bound $|M(x)| = O(x^{1/2+\varepsilon})$ would
    fail, replaced by $|M(x)| = O(x^{\sigma + \varepsilon})$ for
    $\sigma > 1/2$.
\end{enumerate}

None of these affect factoring hardness. The security of RSA does
not depend on the size of $\pi(x) - \text{Li}(x)$; it depends on
the hardness of factoring specific integers.

The nuance: a disproof would mean that primes can cluster or gap
more than expected. In principle, this could affect the output
distribution of prime-generation algorithms (e.g., the probability
that a random odd number in $[2^{1023}, 2^{1024}]$ is prime). But
in practice, the error term is already negligible: for 1024-bit
primes, the density is $\sim 1/710$ regardless of RH, and the
fluctuation term is of order $2^{512}/710^2 \approx 10^{-6}$---far
too small to exploit.

\begin{myth}[``The proof technique itself (Parseval Bridge, Mellin
transform) could be weaponized.'']
\label{myth:technique}
\textbf{False.}
\end{myth}

The Parseval Bridge converts an $L^2$ norm on $(0,1)$ to an $L^2$
norm on the critical line. The Mellin transform connects time-domain
and frequency-domain representations. These are standard tools in
harmonic analysis and signal processing. They are already used in
cryptanalysis (e.g., the Number Field Sieve uses lattice reduction,
which involves Fourier analysis). The Cathedral does not introduce
any new algorithmic technique---it provides a new \emph{structural
understanding} of a known equivalence.

Structural understanding is valuable for mathematics but not
directly useful for algorithm design. Knowing \emph{why} $d_N^2 \to 0$
is equivalent to RH does not tell you \emph{how} to factor integers.

% ================================================================
\section{The Precise Separation: Distribution vs. Computation}
\label{sec:separation}
% ================================================================

The fundamental reason that RH does not affect cryptography is a
gap between two types of mathematical knowledge:

\begin{definition}[Distributional Knowledge]
A statement about the \emph{statistical properties} of a class of
numbers. Example: ``There are $\sim x/\ln x$ primes below $x$.''
\end{definition}

\begin{definition}[Computational Knowledge]
A statement about the \emph{algorithmic difficulty} of a specific
task on a given input. Example: ``Factoring $N = PQ$ requires
$\exp(\Omega(n^{1/3}))$ time for $n$-bit $N$.''
\end{definition}

The Riemann Hypothesis is purely distributional:

\begin{center}
\begin{tabular}{@{}ll@{}}
\toprule
\textbf{What RH tells you} & \textbf{What RH does NOT tell you} \\
\midrule
How many primes $\leq x$ & Which integers are prime \\
Error in prime counting & How to factor $N = PQ$ \\
$|M(x)| \leq C\sqrt{x}\log^2 x$ & $\mu(N)$ for specific $N$ \\
$\zeta(s) \neq 0$ for $\re s > 1/2$ & Discrete log of $g^x \bmod p$ \\
Primes are well-distributed & Private key from public key \\
\bottomrule
\end{tabular}
\end{center}

Cryptographic hardness assumptions (factoring, discrete log, LWE)
are \emph{computational} statements. They assert that certain functions
are hard to invert, not that certain statistical distributions have
particular shapes. RH constrains the shape; cryptography needs the
hardness. These are orthogonal.

\subsection{A Formal Argument}

\begin{theorem}[RH Does Not Imply Factoring]
\label{thm:separation}
There is no known reduction from the Integer Factoring Problem to
any decision problem that is resolved by the Riemann Hypothesis.
That is: no algorithm is known that, given oracle access to an
RH oracle (a function that returns ``true'' for $\RH$), can factor
integers in polynomial time.
\end{theorem}

\begin{proof}[Argument]
An ``RH oracle'' is a single-bit oracle: it returns 1 (if RH is
true) or 0 (if RH is false). A single bit of information cannot
provide a polynomial-time speedup for factoring, since the factoring
problem requires $\Omega(n)$ bits of output (the factor itself has
$n/2$ bits).

More generally: RH is a single mathematical statement. Its truth
value is a fixed constant (either $\top$ or $\bot$). Any algorithm
can be modified to hardcode this constant with at most a constant
factor overhead. If ``assuming RH'' gave a factoring algorithm, then
the unconditional algorithm (which simply tries both branches)
would also be polynomial-time.
\end{proof}

\begin{remark}
This argument generalizes: \emph{no fixed mathematical theorem can
break a cryptosystem}, because the theorem's content can be
hardcoded. Only \emph{algorithmic techniques} can break
cryptosystems, and a theorem statement (however deep) is not an
algorithm.
\end{remark}

% ================================================================
\section{What \emph{Could} (Theoretically) Matter}
\label{sec:could}
% ================================================================

While the Cathedral itself is cryptographically neutral, intellectual
honesty requires addressing the one scenario that \emph{could}
theoretically affect cryptography:

\subsection{The Technique Spillover Hypothesis}

The process of \emph{proving} a deep theorem sometimes produces new
mathematical techniques that have unexpected applications. For example:

\begin{itemize}
  \item Wiles' proof of Fermat's Last Theorem developed the theory of
    Galois representations, which later influenced elliptic curve
    cryptography (both positively and as an analytical tool).
  \item Shor's algorithm for quantum factoring uses the Quantum Fourier
    Transform, which is a general-purpose tool with many applications.
\end{itemize}

If a future proof of RH required developing a fundamentally new
technique for understanding prime factorizations (not just
distributions), that technique \emph{might} have algorithmic
applications. This is speculative and has no evidence base.

\begin{claim}[Cathedral Technique Assessment]
The techniques used in the Cathedral---both active and archived---are:
\begin{enumerate}
  \item Nyman--Beurling $L^2$ density (1950s).
  \item B\'aez-Duarte basis (2003).
  \item Mellin transform and Plancherel theorem (19th century).
  \item Abel summation (17th century).
  \item Montgomery--Vaughan mean value theorem (1973).
  \item Rayleigh--Ritz variational principle (19th century).
  \item Sieve methods: Selberg sieve, Vaughan decomposition (1940s--70s).
  \item Perron's formula for Dirichlet series (1908).
  \item Vasyunin cotangent formula and Dedekind sums (1892--1995).
  \item Robin's inequality and the divisor function (1984).
\end{enumerate}
All of these are well-known, decades-old techniques in harmonic
analysis and analytic number theory. None of them provide any
avenue for algorithmic improvement in factoring, discrete log,
or lattice problems. The Cathedral does not introduce new
mathematical techniques---it assembles existing ones in a new
architecture that is machine-verified.
\end{claim}

\subsection{The GRH-Conditional Theorems}

Several results in theoretical computer science are proved
assuming the Generalized Riemann Hypothesis (GRH), which is
stronger than RH:

\begin{center}
\begin{tabular}{@{}ll@{}}
\toprule
\textbf{GRH-Conditional Result} & \textbf{Impact if GRH Proved} \\
\midrule
Miller's deterministic primality test & Unconditional $O(\log^4 n)$ \\
Polynomial-time class group computation & Unconditional \\
Smallest quadratic non-residue $\leq O(\log^2 p)$ & Unconditional \\
Bach's bounds for generating primes & Tighter constants \\
\bottomrule
\end{tabular}
\end{center}

If GRH were proved, these conditional results would become
unconditional. None of them break any cryptosystem:

\begin{itemize}
  \item \textbf{Miller's test}: Primality testing is already
    polynomial-time unconditionally (AKS, 2002). The GRH-conditional
    Miller test is faster but solves the same problem. Making
    primality testing slightly faster does not help factoring.
  \item \textbf{Class group computation}: Used in computational
    algebraic number theory, not in attacks on deployed systems.
  \item \textbf{Quadratic non-residues}: The smallest quadratic
    non-residue modulo $p$ is already known to be $O(\sqrt{p}\log p)$
    unconditionally (P\'olya--Vinogradov). The GRH improvement to
    $O(\log^2 p)$ does not provide an attack vector.
\end{itemize}

\begin{remark}[GRH Actually Helps Security Proofs]
Several cryptographic security proofs \emph{assume} GRH. If GRH
were proved, these proofs would become unconditional, making the
security guarantees \emph{stronger}, not weaker. Examples include
certain proofs about the Goldwasser--Micali encryption scheme and
the Blum--Blum--Shub pseudorandom generator.
\end{remark}

% ================================================================
\section{Post-Quantum Considerations}
\label{sec:postquantum}
% ================================================================

Modern cryptography is transitioning to post-quantum algorithms
(NIST PQC standards: CRYSTALS-Kyber, CRYSTALS-Dilithium, SPHINCS+).
These are based on:

\begin{itemize}
  \item Lattice problems (Learning with Errors, LWE).
  \item Hash-based signatures.
  \item Code-based cryptography.
\end{itemize}

None of these have any mathematical connection to the Riemann zeta
function, prime distribution, or the Nyman--Beurling criterion. The
hardness of LWE is related to worst-case lattice problems
(Shortest Vector Problem), which live in a completely different
branch of mathematics (geometry of numbers, rather than analytic
number theory).

\begin{claim}[Post-Quantum Immunity]
The Cathedral's result is irrelevant to post-quantum cryptographic
security. The NIST PQC standards are not affected by any resolution
of the Riemann Hypothesis.
\end{claim}

The actual threat to pre-quantum cryptography (RSA, ECC) is
\textbf{Shor's algorithm}, which factors integers and computes
discrete logarithms in quantum polynomial time. Shor's algorithm
has nothing to do with the Riemann Hypothesis---it uses the quantum
Fourier transform to find the period of modular exponentiation, which
is a group-theoretic technique, not a number-theoretic one.

% ================================================================
\section{Archive Audit: 96 Abandoned Approaches}
\label{sec:archive}
% ================================================================

The Cathedral preserves 96 archived Lean files documenting approaches
that were explored and abandoned during the 23-day construction. A
natural question arises: could any of these abandoned techniques
be weaponized, even if the main proof architecture cannot?

We audited all 96 archived files, categorized by technique:

\begin{center}
\small
\begin{tabular}{@{}llcl@{}}
\toprule
\textbf{Category} & \textbf{Technique} & \textbf{Files} & \textbf{Finding} \\
\midrule
Sieve methods & Selberg, Vaughan, parity barrier & 5 & No threat \\
Perron/contour & Contour shifts, Perron kernel & 4 & No threat \\
Spectral theory & Eigenvalues, Hilbert--P\'olya & 3 & No threat \\
Robin's inequality & $\sigma(n) < e^\gamma n\ln\ln n$ & 6 & No threat \\
Vasyunin/cotangent & Dedekind sums, piecewise FTC & 13 & No threat \\
Mellin transforms & Floor Mellin, Dirichlet series & 6 & No threat \\
Exotic approaches & Bessel, winding number, QFT & 18 & No threat \\
Duplicate/structural & Superseded declarations & 41 & No threat \\
\bottomrule
\end{tabular}
\end{center}

\subsection{The Archive as Counter-Evidence}

The archive \emph{strengthens} the security assessment. Each archived
file documents an approach that was explored and \emph{failed}. These
96 failures constitute concrete evidence that the mathematical tools
used in the Cathedral do not produce computational shortcuts.

\begin{claim}[Archive as Counter-Evidence]
If sieve methods, Perron's formula, or Mellin transform techniques
could be repurposed for integer factoring or discrete logarithm
computation, the archived approaches that used these techniques
would have \emph{succeeded}---they would have produced unexpected
algorithmic results instead of mathematical dead ends. Their failure
is evidence \emph{for} the distributional-vs-computational separation
(Section~\ref{sec:separation}).
\end{claim}

The distinction is important: these are not hypothetical arguments
about what these tools \emph{might} do. They are empirical records
of what these tools \emph{actually did} when applied to a
concrete problem. The answer: they bounded $L^2$ norms, computed
Gram matrix entries, and controlled Dirichlet polynomial averages.
At no point did any approach produce information about the
factorization of specific integers.

\subsection{The Rust Experiments (211,202 lines)}

The companion Rust codebase computes Gram matrix entries, eigenvalues,
and Rayleigh quotients using exact rational arithmetic for
$N \leq 5000$. It processes the integers $2, 3, \ldots, N$
using $\gcd$ and $\operatorname{lcm}$ operations. This is publicly
reproducible computation on small public integers. No private keys,
secrets, or cryptographic targets are involved at any point in the
numerical pipeline.

% ================================================================
\section{Symmetric Cryptography and Hash Functions}
\label{sec:symmetric}
% ================================================================

For completeness: symmetric cryptography (AES, ChaCha20) and hash
functions (SHA-2, SHA-3, BLAKE3) have no connection to prime numbers
or number theory whatsoever. Their security is based on:

\begin{itemize}
  \item Confusion and diffusion (Shannon, 1949).
  \item Resistance to linear and differential cryptanalysis.
  \item The random oracle model (for hash functions).
\end{itemize}

The Riemann Hypothesis could be proved, disproved, or shown to be
independent of ZFC, and AES-256 would remain equally secure (or
insecure) as it is today.

% ================================================================
\section{What a Responsible Disclosure Looks Like}
\label{sec:disclosure}
% ================================================================

Since the Cathedral is cryptographically neutral, no security
advisory is warranted. However, for the benefit of responsible
disclosure processes, we state explicitly:

\begin{enumerate}
  \item \textbf{No CVE is needed.} The Cathedral does not disclose
    a vulnerability in any system.
  \item \textbf{No key material is compromised.} The proof involves
    no private keys, secrets, or confidential data.
  \item \textbf{No algorithm is accelerated.} The proof does not
    provide a faster factoring, discrete log, or lattice algorithm.
  \item \textbf{No system requires patching.} No software update
    is needed in response to this work.
  \item \textbf{No embargo is appropriate.} The work is pure
    mathematics and should be published openly.
\end{enumerate}

% ================================================================
\section{A Note for Journalists}
\label{sec:journalists}
% ================================================================

If you are writing about the Cathedral for a general audience,
here are accurate and inaccurate framings:

\begin{center}
\begin{tabular}{@{}p{0.45\textwidth}p{0.45\textwidth}@{}}
\toprule
\textbf{\color{green!50!black}Accurate} &
\textbf{\color{red!70!black}Inaccurate} \\
\midrule
``Researchers built a computer-verified proof about the oldest
unsolved problem in math.'' &
``Researchers cracked the code behind internet encryption.'' \\[0.5em]
``The result translates the Riemann Hypothesis into an equivalent
form, but does not solve it.'' &
``The Riemann Hypothesis has been solved / nearly solved.'' \\[0.5em]
``This work has no impact on cryptographic security.'' &
``Your passwords / bank accounts may be at risk.'' \\[0.5em]
``The proof is about the statistical distribution of prime numbers.'' &
``The proof reveals a pattern in primes that hackers could exploit.'' \\[0.5em]
``If anything, proving RH would strengthen crypto security proofs.'' &
``Proving RH would break the internet.'' \\
\bottomrule
\end{tabular}
\end{center}

% ================================================================
\section{Conclusion}
\label{sec:conclusion}
% ================================================================

The Cathedral is a machine-verified proof that the Riemann Hypothesis
is equivalent to the decay of the Nyman--Beurling distance. It is a
contribution to pure mathematics and proof engineering, with no
direct, indirect, or foreseeable implications for cryptographic
security.

The disconnect between RH and cryptography is not accidental---it
reflects a deep structural separation between \emph{distributional}
properties of primes (which RH constrains) and \emph{computational}
hardness assumptions (which cryptography requires). These are
different mathematical universes. Progress in one does not translate
to the other.

If the Riemann Hypothesis is eventually proved:
\begin{itemize}
  \item Mathematicians will celebrate.
  \item Cryptographers will note that several conditional security
    proofs are now unconditional (a mild strengthening).
  \item Security practitioners will change nothing.
  \item Your bank account will be fine.
\end{itemize}

If the Riemann Hypothesis is eventually disproved:
\begin{itemize}
  \item Mathematicians will be astonished.
  \item Cryptographers will note that some error bounds in prime
    generation are slightly looser than assumed (negligible in
    practice).
  \item Security practitioners will \emph{still} change nothing.
  \item Your bank account will \emph{still} be fine.
\end{itemize}

The primes are not the lock. The primes are the landscape.
Knowing the landscape perfectly does not give you the key.

\vspace{0.5cm}
\begin{flushright}
\emph{Sleep well. Your encryption is safe.}
\end{flushright}

\vspace{1cm}

% ================================================================
% BIBLIOGRAPHY
% ================================================================
\begin{thebibliography}{99}

\bibitem{rsa}
R.~Rivest, A.~Shamir, and L.~Adleman, \emph{A method for obtaining
digital signatures and public-key cryptosystems}, Comm. ACM
\textbf{21} (1978), 120--126.

\bibitem{aks}
M.~Agrawal, N.~Kayal, and N.~Saxena, \emph{PRIMES is in P},
Annals of Mathematics \textbf{160} (2004), 781--793.

\bibitem{miller}
G.L.~Miller, \emph{Riemann's hypothesis and tests for primality},
J. Comput. System Sci. \textbf{13} (1976), 300--317.

\bibitem{shor}
P.~Shor, \emph{Polynomial-time algorithms for prime factorization
and discrete logarithms on a quantum computer}, SIAM J. Comput.
\textbf{26} (1997), 1484--1509.

\bibitem{nist-pqc}
NIST, \emph{Post-Quantum Cryptography Standardization},
\url{https://csrc.nist.gov/projects/post-quantum-cryptography}, 2024.

\bibitem{montgomery1973}
H.L.~Montgomery, \emph{The pair correlation of zeros of the zeta
function}, Proc. Symp. Pure Math. \textbf{24} (1973), 181--193.

\bibitem{baezduarte2003}
L.~B\'aez-Duarte, \emph{A strengthening of the Nyman--Beurling
criterion}, Atti Accad. Naz. Lincei Rend. Mat. Appl. \textbf{14}
(2003), 5--11.

\bibitem{lwe}
O.~Regev, \emph{On lattices, learning with errors, random linear
codes, and cryptography}, J. ACM \textbf{56} (2009), 1--40.

\end{thebibliography}

\end{document}
